\section{Cel}

Ten dokument testuje darmowy workflow:

\begin{quote}
GitHub Actions kompiluje przykładowy projekt LaTeX do PDF bez VPS i bez Rendera.
\end{quote}

\section{Zakres MVP}

\begin{center}
\begin{tabular}{ll}
\toprule
Funkcja & Status \\
\midrule
Upload TEX lokalnie & OK \\
Upload ZIP lokalnie & OK \\
Slayer Paper Mode & OK \\
Historia buildów lokalnie & OK \\
GitHub Actions PDF builder & Demo \\
\bottomrule
\end{tabular}
\end{center}

\section{Wniosek}

To nie zastępuje backendu online, ale pozwala publicznie pokazać, że repo potrafi automatycznie budować PDF z projektu LaTeX.

\section{Następny krok}

Docelowo frontend z GitHub Pages może wysyłać pliki do osobnego backendu na VPS albo do środowiska udostępnionego przez Slayer Project.
